% Les remerciements.
%
% AUCUN NOM DE PERSONNE EN CLAIR. L'encadrant de stage est nommé par son marqueur, comme sur la page
% de garde, et le marqueur reste vide tant que `noms.tex` n'est pas déposé. La phrase se compose
% alors sans nom, ce qui est le comportement attendu en état de brouillon.
%
% Rien n'est affirmé qui ne soit établi. Le seul fait dont ce dossier porte la trace est que
% l'encadrant a remis le cahier de charges et laissé toute liberté d'action dans le cadre qu'il
% fixe : c'est ce que les remerciements disent, et rien de plus.

\chapter*{Remerciements}
\addcontentsline{toc}{chapter}{Remerciements}

Je remercie mon encadrant de
stage\siRenseigne{\marqueurEncadrantProfessionnel}{, \marqueurEncadrantProfessionnel,} pour la
confiance qu'il m'a accordée. Il m'a remis le cahier de charges, en a précisé le contenu, et m'a
laissé toute liberté sur la démarche à suivre dans le cadre qu'il avait fixé. Cette liberté a
déplacé vers moi la responsabilité du cadrage, et c'est ce qui a rendu ce travail formateur.

Je remercie le service d'accueil et d'admission qui a hébergé ce travail, à
l'\marqueurOrganismeAccueil{}, pour son accueil et pour le temps qu'il m'a consacré au poste
d'observation.

Je remercie enfin\siRenseigne{\marqueurEtablissementFormation}{ l'\marqueurEtablissementFormation{},}
pour la formation qui m'a permis d'aborder ce travail.
